\documentclass[11pt]{article}
\usepackage[T1]{fontenc}
\usepackage[utf8]{inputenc}
\usepackage{lmodern}
\usepackage{geometry}
\usepackage{hyperref}
\usepackage{enumitem}
\usepackage{verbatim}
\geometry{margin=1in}

\title{Detailed Implementation Plan for a Hybrid Exasim Architecture}
\author{}
\date{}

\begin{document}
\maketitle

\section*{1. Summary}

Yes. The maintainable implementation strategy is to \textbf{shrink the template boundary to the model/kernel edge} and move everything else into a compiled core.

A concrete plan for Exasim is:

\begin{itemize}[leftmargin=*]
\item keep \texttt{MyModel} and the model-dependent Kokkos kernels templated
\item replace \texttt{CDiscretization<M>}, \texttt{CSolution<M>}, \texttt{CSolver<M>}, \texttt{CPreconditioner<M>} with non-templated core classes
\item introduce a small host-side \texttt{ModelOps} / \texttt{ModelContext} adapter that launches model-specialized kernels
\item keep \texttt{ModelContext} as a simple struct, not a heavy class
\item pass \texttt{ModelContext} into the \texttt{CSolution(...)} and \texttt{CDiscretization(...)} constructors so constructor-time initialization can use model-dependent init operations
\item compile the backend into \texttt{libexasim\_core.a}
\item keep \texttt{run<M>} and \texttt{ExasimSolver<M>} as thin templated frontends over the non-templated core
\item isolate the old ABI path behind one compatibility adapter instead of keeping \texttt{AbiAdapter} in the main architecture
\end{itemize}

\section*{2. Assumptions}

\begin{itemize}[leftmargin=*]
\item The current template spread is the one visible in:
\begin{itemize}[leftmargin=*]
\item \texttt{backend/Discretization/discretization.h}
\item \texttt{backend/Solution/solution.h}
\item \texttt{backend/Solver/solver.h}
\item \texttt{backend/Preconditioning/preconditioner.h}
\end{itemize}
\item The current public package is the header-only \texttt{Exasim::headers} target in \texttt{install/CMakeLists.txt}.
\item Model-dependent launches already have a natural boundary in:
\begin{itemize}[leftmargin=*]
\item \texttt{include/exasim/drivers.hpp}
\item \texttt{include/exasim/kernels/}
\end{itemize}
\item \texttt{readsolstruct} in \texttt{backend/Discretization/readbinaryfiles.hpp} computes initial fields through \texttt{cpuInit*Driver} calls, so if \texttt{cpuInit*} is moved into \texttt{ModelOps}, \texttt{ModelContext} must be available during construction.
\item Host-side dispatch through function pointers is acceptable if it happens at driver-launch granularity, not inside inner device loops.
\end{itemize}

\section*{3. Design reasoning}

The current \texttt{library-port} over-templates the backend.

The symptom is that \texttt{M} leaks into:

\begin{itemize}[leftmargin=*]
\item discretization
\item solution
\item solver
\item preconditioner
\item point locator
\item output/QoI paths
\item install/public headers
\end{itemize}

That causes:

\begin{itemize}[leftmargin=*]
\item template contagion
\item backend \texttt{.hpp} growth
\item \texttt{.cpp} files being included into headers
\item \texttt{inline}/ODR workarounds like \texttt{backend/Model/KokkosDrivers.cpp}
\item a very large public transitive include graph
\end{itemize}

The preferred lifecycle is also important:

\begin{itemize}[leftmargin=*]
\item \texttt{ModelContext} is created at the templated API boundary
\item \texttt{CSolution(...)} receives \texttt{ModelContext} at construction time
\item \texttt{CSolution(...)} passes \texttt{ModelContext} into \texttt{CDiscretization(...)} at construction time
\item constructor-time paths such as \texttt{readsolstruct} can immediately use \texttt{getModelContext()}
\item \texttt{bindModelContext(...)} is no longer the primary initialization path; at most it is a secondary/rebind mechanism
\end{itemize}

The better boundary is:

\begin{itemize}[leftmargin=*]
\item \textbf{templated:} only code that truly needs \texttt{M::flux}, \texttt{M::source}, \texttt{M::fbou\_hdg}, Jacobians, and compile-time sizes
\item \textbf{non-templated:} orchestration, ownership, file I/O, solver flow, mesh/preprocessing, point location, output plumbing
\end{itemize}

That means the current \texttt{EXASIM\_DRIVER\_CALL(...)} model should be replaced by a runtime \texttt{ModelOps} object whose function pointers point to \textbf{host-side wrappers} specialized on \texttt{M}. Those wrappers then launch the templated kernels.

\section*{4. Proposed implementation}

\subsection*{Phase 0: Freeze invariants and scope}

\textbf{Affected files:}
\begin{itemize}[leftmargin=*]
\item \texttt{docs/library\_port\_implementation\_strategy.tex}
\item add a new architecture note under \texttt{docs/} or \texttt{docs/03-internals/}
\end{itemize}

\textbf{Plan:}
\begin{itemize}[leftmargin=*]
\item Define the non-negotiable numerical and API invariants before refactoring.
\item Lock down that HDG assembly structure must remain unchanged.
\item Write a one-page boundary note: ``templated kernel edge / compiled runtime core.''
\item Identify which public entry points must remain source-compatible:
\begin{itemize}[leftmargin=*]
\item \texttt{exasim::run<M>(...)}
\item \texttt{ExasimSolver<M>}
\item legacy \texttt{backend/Main/main.cpp}
\end{itemize}
\end{itemize}

\textbf{Deliverable:}
\begin{itemize}[leftmargin=*]
\item design note with exact boundary and migration rules
\end{itemize}

\subsection*{Phase 1: Introduce a runtime model interface}

\textbf{New files:}
\begin{itemize}[leftmargin=*]
\item \texttt{include/exasim/model\_ops.hpp}
\item \texttt{include/exasim/model\_context.hpp}
\item \texttt{include/exasim/model\_adapter.hpp}
\item \texttt{backend/Model/model\_ops.cpp} only if shared non-template helpers are needed
\end{itemize}

\textbf{Plan:}
\begin{itemize}[leftmargin=*]
\item Define a \texttt{ModelOps} struct of host-callable function pointers for each driver family currently reached through \texttt{EXASIM\_DRIVER\_CALL(...)}.
\item Split the ops by category:
\begin{itemize}[leftmargin=*]
\item volume: \texttt{flux}, \texttt{source}, \texttt{sourcew}, \texttt{eos}, \texttt{tdfunc}
\item boundary/interface: \texttt{fbou}, \texttt{fhat}, \texttt{fint}, \texttt{fext}, \texttt{ubou}
\item init: \texttt{initu}, \texttt{initudg}, \texttt{initwdg}, \texttt{initodg}
\item output/QoI/vis: \texttt{output}, \texttt{monitor}, \texttt{qoi\_volume}, \texttt{qoi\_boundary}, \texttt{vis\_scalars}, \texttt{vis\_vectors}, \texttt{vis\_tensors}
\end{itemize}
\item Define \texttt{ModelContext} as:
\begin{itemize}[leftmargin=*]
\item \texttt{const ModelOps* ops}
\item model metadata: \texttt{nd}, \texttt{ncu}, \texttt{ncw}, \texttt{nco}, \texttt{disc}, maybe feature flags
\item optional opaque pointer for model-owned state, if ever needed later
\end{itemize}
\item Keep \texttt{ModelContext} as a plain struct used only to carry dispatch pointers and model metadata. Do not make it a heavy class unless it later needs ownership or invariant logic.
\item Build \texttt{ModelContext} at the templated boundary only, then pass it into \texttt{CSolution(...)} and \texttt{CDiscretization(...)} as part of normal construction.
\item \texttt{cpuInitu}, \texttt{cpuInitudg}, \texttt{cpuInitwdg}, and \texttt{cpuInitodg} must be represented in \texttt{ModelOps}, because initialization is model-dependent and may run during \texttt{readsolstruct}.
\end{itemize}

\textbf{Important implementation detail:}
\begin{itemize}[leftmargin=*]
\item These function pointers must point to host functions that launch Kokkos kernels specialized on \texttt{M}.
\item Do not use runtime polymorphism inside device code.
\end{itemize}

\textbf{Deliverable:}
\begin{itemize}[leftmargin=*]
\item a runtime model interface that can replace the macro-based dual dispatch
\end{itemize}

\subsection*{Phase 2: Build templated model adapters}

\textbf{New files:}
\begin{itemize}[leftmargin=*]
\item \texttt{include/exasim/model\_adapter\_impl.hpp}
\item \texttt{backend/ModelAdapters/abi\_adapter.cpp}
\item generated or handwritten adapter TUs, for example:
\begin{itemize}[leftmargin=*]
\item \texttt{apps/library\_example/poisson2d/model\_adapter.cpp}
\item generated codegen output \texttt{.../my\_model\_adapter.cpp}
\end{itemize}
\end{itemize}

\textbf{Plan:}
\begin{itemize}[leftmargin=*]
\item Implement \texttt{template <class M> ModelOps make\_model\_ops();}
\item Each function in \texttt{ModelOps} forwards to one current templated driver in \texttt{include/exasim/drivers.hpp}.
\item Example pattern:
\begin{itemize}[leftmargin=*]
\item \texttt{ops.flux = \&flux\_driver\_adapter<M>;}
\item where \texttt{flux\_driver\_adapter<M>(...)} calls \texttt{exasim::FluxDriver<M>(...)}
\end{itemize}
\item This is the key move that narrows the template boundary without losing specialization.
\end{itemize}

\textbf{Deliverable:}
\begin{itemize}[leftmargin=*]
\item one model adapter TU per model family or executable
\item no need for large backend classes to depend on \texttt{M}
\end{itemize}

\subsection*{Phase 3: Replace templated backend classes with core classes}

\textbf{Affected files:}
\begin{itemize}[leftmargin=*]
\item \texttt{backend/Discretization/discretization.h}
\item \texttt{backend/Discretization/discretization.hpp}
\item \texttt{backend/Solution/solution.h}
\item \texttt{backend/Solution/solution.hpp}
\item \texttt{backend/Solver/solver.h}
\item \texttt{backend/Solver/solver.hpp}
\item \texttt{backend/Preconditioning/preconditioner.h}
\item \texttt{backend/Preconditioning/preconditioner.hpp}
\item \texttt{backend/PointLocator/pointlocator.h}
\item \texttt{backend/PointLocator/pointlocator.hpp}
\end{itemize}

\textbf{Plan:}
\begin{itemize}[leftmargin=*]
\item Introduce non-templated replacements:
\begin{itemize}[leftmargin=*]
\item \texttt{CDiscretization}
\item \texttt{CSolution}
\item \texttt{CSolver}
\item \texttt{CPreconditioner}
\end{itemize}
\item Replace \texttt{CDiscretization<M>\&} references with \texttt{CDiscretization\&}.
\item Store \texttt{const ModelContext*} or \texttt{const ModelOps*} inside \texttt{CDiscretization}, and initialize it during construction.
\item Update constructors so \texttt{ModelContext} is part of the construction dependency graph:
\begin{itemize}[leftmargin=*]
\item \texttt{CDiscretization(const ModelContext\&, ...)}
\item \texttt{CSolution(const ModelContext\&, ...)}
\item \texttt{CSolver(...)} and \texttt{CPreconditioner(...)} receive it directly only if they need constructor-time model access
\end{itemize}
\item Keep optional bind methods only as secondary/rebind hooks:
\begin{itemize}[leftmargin=*]
\item \texttt{CDiscretization::bindModelContext(const ModelContext\&)}
\item \texttt{CSolver::bindModelContext(const ModelContext\&)} if solver-side access is needed
\item \texttt{CPreconditioner::bindModelContext(const ModelContext\&)} if preconditioner-side access is needed
\item \texttt{CSolution::bindModelContext(const ModelContext\&)} as the top-level coordinator for rare rebind flows
\end{itemize}
\item Keep the constructor flow centered on \texttt{CSolution(...)}. The constructor now receives \texttt{ModelContext} and propagates it to \texttt{disc} immediately so constructor-time paths such as \texttt{readsolstruct} can use model-dependent init operations.
\item If retained, \texttt{CSolution::bindModelContext(...)} is not the main public initialization path. It is only for controlled rebind or reconstruction scenarios.
\item Add accessor functions instead of exposing raw members:
\begin{itemize}[leftmargin=*]
\item \texttt{CDiscretization::getModelContext()}
\item \texttt{CSolution::getModelContext()} if the solution object needs to expose the bound context
\end{itemize}
\item Use \texttt{getModelContext()} as the checked accessor name rather than \texttt{model\_context()} or \texttt{ModelContext()} to avoid confusion with the \texttt{ModelContext} type itself.
\item Add an internal validity check on bind:
\begin{itemize}[leftmargin=*]
\item \texttt{disc.common.nd == model.nd}
\item \texttt{disc.common.ncu == model.ncu}
\item \texttt{disc.common.ncw == model.ncw}
\item discretization kind is compatible with \texttt{model.disc}
\end{itemize}
\item Convert all current template-heavy helper signatures in backend \texttt{.hpp} files to operate on the core types plus \texttt{ModelContext}.
\end{itemize}

\textbf{Deliverable:}
\begin{itemize}[leftmargin=*]
\item non-templated backend class graph
\item template parameter removed from the core ownership and orchestration classes
\item constructor-time \texttt{ModelContext} availability for all initialization paths that need model-dependent operations
\item checked access through \texttt{getModelContext()} rather than raw direct access
\end{itemize}

\subsection*{Phase 4: Eliminate \texttt{EXASIM\_DRIVER\_CALL(...)} from the main backend}

\textbf{Affected files:}
\begin{itemize}[leftmargin=*]
\item \texttt{include/exasim/detail/driver\_dispatch.hpp}
\item all backend call sites using \texttt{EXASIM\_DRIVER\_CALL}, especially:
\begin{itemize}[leftmargin=*]
\item \texttt{backend/Discretization/residual.hpp}
\item \texttt{backend/Discretization/uequation.hpp}
\item \texttt{backend/Discretization/uresidual.hpp}
\item \texttt{backend/Discretization/qoicalculation.hpp}
\item \texttt{backend/Discretization/getuhat.hpp}
\item \texttt{backend/Discretization/discretization.hpp}
\item \texttt{backend/Discretization/postdiscretization.hpp}
\end{itemize}
\end{itemize}

\textbf{Plan:}
\begin{itemize}[leftmargin=*]
\item Replace macro dispatch with explicit runtime calls:
\begin{itemize}[leftmargin=*]
\item \texttt{disc.getModelContext().ops->flux(...)}
\item \texttt{disc.getModelContext().ops->source(...)}
\item etc.
\end{itemize}
\item Keep the old macro only inside the compatibility layer during migration.
\item Make backend residual/assembly code ordinary non-template code unless the algorithm truly depends on compile-time dimensions.
\item Route constructor-time initialization through the same runtime dispatch system, for example \texttt{getModelContext().ops->initu(...)} and related init operations instead of direct legacy-only \texttt{cpuInit*Driver} calls.
\end{itemize}

\textbf{Deliverable:}
\begin{itemize}[leftmargin=*]
\item macro removed from the main backend path
\item much smaller template footprint in \texttt{backend/Discretization/*.hpp}
\end{itemize}

\subsection*{Phase 5: Remove \texttt{.cpp} inclusion from headers and build a static core}

\textbf{Affected files:}
\begin{itemize}[leftmargin=*]
\item \texttt{backend/Discretization/discretization.hpp}
\item \texttt{backend/Discretization/postdiscretization.hpp}
\item \texttt{backend/Model/KokkosDrivers.cpp}
\item \texttt{backend/Model/ModelDrivers.cpp}
\item \texttt{install/CMakeLists.txt}
\end{itemize}

\textbf{Plan:}
\begin{itemize}[leftmargin=*]
\item Stop \texttt{\#include}-ing \texttt{backend/Model/KokkosDrivers.cpp} and similar \texttt{.cpp} files from backend headers.
\item Build a real static target, for example:
\begin{itemize}[leftmargin=*]
\item \texttt{add\_library(exasim\_core STATIC ...)}
\item export as \texttt{Exasim::core}
\end{itemize}
\item Move non-template implementations into normal \texttt{.cpp} translation units.
\item At this point, \texttt{inline} on those free functions should no longer be necessary.
\end{itemize}

\textbf{Deliverable:}
\begin{itemize}[leftmargin=*]
\item \texttt{libexasim\_core.a}
\item no more ``included \texttt{.cpp} as implementation header'' pattern in the core
\end{itemize}

\subsection*{Phase 6: Refactor public API into thin templated wrappers}

\textbf{Affected files:}
\begin{itemize}[leftmargin=*]
\item \texttt{include/exasim/run.hpp}
\item \texttt{include/exasim/solver\_facade.hpp}
\item public wrappers under \texttt{include/exasim/}
\end{itemize}

\textbf{Plan:}
\begin{itemize}[leftmargin=*]
\item Keep \texttt{run<M>} templated, but make it thin:
\begin{itemize}[leftmargin=*]
\item build \texttt{ModelContext ctx = make\_model\_context<M>();}
\item construct \texttt{CSolution(ctx, ...)} through the usual file/runtime path
\item then call non-templated \texttt{run\_impl(...);}
\end{itemize}
\item Keep \texttt{ExasimSolver<M>} templated only at construction/binding points:
\begin{itemize}[leftmargin=*]
\item \texttt{ExasimSolver<M>} owns \texttt{std::unique\_ptr<CSolution>}
\item it constructs the solution object with a ready \texttt{ModelContext}
\end{itemize}
\item Remove backend-heavy includes from public headers where possible.
\item Make public headers depend on stable declarations, not large backend internals.
\end{itemize}

\textbf{Deliverable:}
\begin{itemize}[leftmargin=*]
\item same user-facing syntax
\item much thinner template layer
\end{itemize}

\subsection*{Phase 7: Isolate the legacy ABI path}

\textbf{Affected files:}
\begin{itemize}[leftmargin=*]
\item \texttt{include/exasim/detail/abi\_adapter.hpp}
\item \texttt{backend/Main/main.cpp}
\item \texttt{backend/Model/libpdemodel.hpp}
\item \texttt{backend/Model/libpdemodel.cpp}
\item legacy \texttt{ModelDrivers.cpp} / \texttt{KokkosDrivers.cpp}
\end{itemize}

\textbf{Plan:}
\begin{itemize}[leftmargin=*]
\item Replace ``ABI is one branch of the main architecture'' with ``ABI is one adapter implementation of \texttt{ModelOps}.''
\item Create \texttt{make\_model\_ops\_abi()} that binds the old driver symbols.
\item Keep \texttt{backend/Main/main.cpp} using that adapter.
\item Remove ABI branching from the main template/runtime design.
\end{itemize}

\textbf{Deliverable:}
\begin{itemize}[leftmargin=*]
\item ABI path survives
\item but no longer shapes the public or core architecture
\end{itemize}

\subsection*{Phase 8: Packaging and install model}

\textbf{Affected files:}
\begin{itemize}[leftmargin=*]
\item \texttt{install/CMakeLists.txt}
\item \texttt{tests/install\_consumer/CMakeLists.txt}
\item \texttt{tests/multi\_tu\_consumer/}
\end{itemize}

\textbf{Plan:}
\begin{itemize}[leftmargin=*]
\item Replace header-only export with:
\begin{itemize}[leftmargin=*]
\item \texttt{Exasim::core} static library
\item optional small \texttt{Exasim::model\_api} interface target if needed
\end{itemize}
\item Install only the public API headers and genuinely templated kernel headers.
\item Stop installing large parts of \texttt{backend/} as public transitive implementation detail.
\item Keep consumer tests:
\begin{itemize}[leftmargin=*]
\item out-of-tree install test
\item multi-TU consumer test
\item one handwritten model consumer
\item one generated model consumer
\end{itemize}
\end{itemize}

\textbf{Deliverable:}
\begin{itemize}[leftmargin=*]
\item cleaner install tree
\item smaller public header surface
\item simpler consumer story
\end{itemize}

\subsection*{Phase 9: Controlled migration sequence}

\textbf{Order of execution:}
\begin{enumerate}[leftmargin=*]
\item add \texttt{ModelOps} and \texttt{ModelContext}
\item add templated adapters for current kernels
\item introduce core classes in parallel with current templated classes
\item migrate one path end-to-end:
\begin{itemize}[leftmargin=*]
\item suggested pilot: handwritten \texttt{poisson2d}
\end{itemize}
\item build \texttt{Exasim::core}
\item convert \texttt{run<M>} and \texttt{ExasimSolver<M>}
\item migrate ABI path behind adapter
\item delete old wide-template path after equivalence is proven
\end{enumerate}

This sequencing matters. Do not flip the whole repo at once.

\section*{5. Risks}

\begin{itemize}[leftmargin=*]
\item The biggest technical risk is placing too much behind runtime dispatch and losing compile-time specialization where HDG kernels need it.
\item The second risk is trying to make the core non-templated too quickly; some helper layers in \texttt{backend/Discretization/} and \texttt{backend/Solver/} may still carry hidden compile-time assumptions.
\item The migration can temporarily increase complexity if old and new paths coexist for too long.
\item Constructor-time \texttt{ModelContext} means the core constructors depend on the model-dispatch setup earlier than before; this is correct for Exasim if initialization uses \texttt{cpuInit*} through \texttt{ModelOps}, but it raises the bar for constructor correctness.
\item GPU builds are the most sensitive phase because kernel launch wrappers, device compilation, and linkage all change at once.
\item The current branch also has local worktree changes, so any implementation should start from a clean branch.
\end{itemize}

\section*{6. Verification plan}

For each phase, require these checks.

\textbf{Numerical equivalence}
\begin{itemize}[leftmargin=*]
\item same residual/Jacobian behavior for one HDG scalar case
\item same QoI and output for one steady case
\item same time-dependent behavior for one DIRK case
\end{itemize}

\textbf{Build/link correctness}
\begin{itemize}[leftmargin=*]
\item single-TU consumer
\item multi-TU consumer
\item install-tree \texttt{find\_package(Exasim)} consumer
\item no \texttt{.cpp} implementation includes in public/core headers after Phase 5
\item verify that constructor-time \texttt{readsolstruct} can compute initial fields through \texttt{ModelOps} init functions
\end{itemize}

\textbf{Backend coverage}
\begin{itemize}[leftmargin=*]
\item CPU single-rank
\item CPU MPI
\item NVIDIA GPU
\item one ABI-legacy executable still working through the compatibility adapter
\end{itemize}

\textbf{Architecture acceptance criteria}
\begin{itemize}[leftmargin=*]
\item \texttt{backend/Main/main.cpp} no longer determines the main design
\item \texttt{AbiAdapter} no longer appears in core class templates
\item \texttt{KokkosDrivers.cpp} is compiled, not included
\item \texttt{run<M>} and \texttt{ExasimSolver<M>} are thin wrappers, not owners of the backend architecture
\item \texttt{install/CMakeLists.txt} exports \texttt{Exasim::core}, not only \texttt{Exasim::headers}
\item \texttt{CSolution(...)} and \texttt{CDiscretization(...)} receive \texttt{ModelContext} at construction time
\item \texttt{getModelContext()} is the standard checked accessor for bound model dispatch
\end{itemize}

If you want, I can turn this into a \textbf{work-package plan with estimated effort, file ownership, and merge order}, or into a \textbf{refactor RFC document} under \texttt{docs/}.

\end{document}
